\section{РЕАЛІЗАЦІЯ, ТЕСТУВАННЯ ТА ЕКСПЛУАТАЦІЯ ПРОГРАМНОЇ СИСТЕМИ}

\subsection{Загальна структура програмної реалізації}

Програмну систему реалізовано мовою Python з використанням Flask для web-рівня, PostgreSQL для збереження даних, Alembic для керування міграціями бази даних та Hugging Face Transformers для запуску локальної мовної моделі. Вихідний код розділено на окремі модулі відповідно до їхньої відповідальності.

Головним файлом запуску є \texttt{app.py}. Він аналізує параметри командного рядка і запускає один із двох режимів роботи:
\begin{itemize}
  \item web-режим, у якому створюється Flask-застосунок;
  \item CLI-режим, який активується параметром \texttt{-t} або \texttt{--terminal}.
\end{itemize}

Основні каталоги програмної системи наведено у таблиці~\ref{tab:project-structure}.

\small
\begin{longtable}{|p{4.2cm}|p{10.2cm}|}
\caption{Структура програмної реалізації}
\label{tab:project-structure}\\
\hline
\textbf{Каталог або файл} & \textbf{Призначення} \\
\hline
\endfirsthead
\hline
\textbf{Каталог або файл} & \textbf{Призначення} \\
\hline
\endhead
\texttt{app.py} & Головна точка входу, вибір web- або CLI-режиму. \\
\hline
\texttt{src/core} & Доменні моделі та менеджер чатів. \\
\hline
\texttt{src/web} & Flask routes, HTML-шаблони, CSS і JavaScript web-інтерфейсу. \\
\hline
\texttt{src/cli} & Реалізація консольного режиму. \\
\hline
\texttt{src/storage} & Абстракція сховища та PostgreSQL-реалізація. \\
\hline
\texttt{src/llm} & Інтерфейс LLM-сервісу, runtime-керування моделлю та Transformers-backend. \\
\hline
\texttt{migrations} & Alembic-міграції схеми бази даних. \\
\hline
\texttt{scripts} & Допоміжні скрипти для завантаження моделей, адаптерів і експериментального fine-tuning. \\
\hline
\texttt{tests} & Автоматизовані тести системи. \\
\hline
\end{longtable}
\normalsize

Такий поділ дозволяє відокремити інтерфейс користувача від бізнес-логіки, сховища даних та LLM-рівня. Це спрощує тестування, підтримку та подальший розвиток системи.

\subsection{Реалізація доменної логіки}

Доменний рівень системи реалізовано у модулі \texttt{src/core}. Основними сутностями є \texttt{Chat}, \texttt{Message} та \texttt{SearchResult}. Вони описують чат, окреме повідомлення та результат пошуку відповідно.

Клас \texttt{ChatManager} є центральним сервісом прикладної логіки. Він не залежить безпосередньо від Flask, PostgreSQL або Transformers, а працює через абстракції \texttt{BaseStorage} та \texttt{BaseLLMService}. Завдяки цьому один і той самий менеджер використовується як у web-інтерфейсі, так і в CLI-режимі.

Основні можливості \texttt{ChatManager}:
\begin{itemize}
  \item створення чатів з унікальними назвами;
  \item перевірка назв чатів на порожнє значення, довжину та дублювання;
  \item отримання списку чатів;
  \item додавання повідомлень користувача й асистента;
  \item генерація відповіді через підключений LLM-сервіс;
  \item streaming-генерація з подальшим збереженням фінальної відповіді;
  \item пошук повідомлень;
  \item перейменування та видалення чатів.
\end{itemize}

При надсиланні повідомлення \texttt{ChatManager} спочатку зберігає повідомлення користувача, формує історію діалогу, передає її до LLM-сервісу, отримує відповідь і зберігає її як повідомлення асистента. Такий підхід забезпечує цілісність історії діалогу та дозволяє моделі враховувати контекст поточного чату.

\subsection{Застосовані принципи проєктування та патерни}

Під час реалізації системи було використано об'єктно-орієнтований підхід з розділенням відповідальності між окремими класами та модулями. Доменні сутності \texttt{Chat}, \texttt{Message} і \texttt{SearchResult} описують основні об'єкти предметної області, а сервіс \texttt{ChatManager} інкапсулює бізнес-логіку роботи з діалогами.

Основні принципи проєктування, використані у системі:
\begin{itemize}
  \item \textbf{Single Responsibility}. Кожен основний компонент має окрему зону відповідальності: \texttt{ChatManager} керує сценаріями чатів, \texttt{PostgresStorage} відповідає за збереження даних, \texttt{TransformersLLMService} інкапсулює роботу з локальною моделлю, а \texttt{LLMRuntime} керує станом завантаженої моделі.
  \item \textbf{Dependency Inversion}. Доменна логіка залежить не від конкретної PostgreSQL-реалізації або Transformers-backend, а від абстракцій \texttt{BaseStorage} і \texttt{BaseLLMService}. Завдяки цьому у тестах можна підставляти \texttt{InMemoryStorage} і \texttt{FakeLLMService}.
  \item \textbf{Open/Closed}. Система може бути розширена новими storage-backend, LLM-backend або конфігураціями моделей без переписування доменної логіки.
  \item \textbf{Separation of Concerns}. Web-рівень, CLI-рівень, бізнес-логіка, сховище даних, LLM-рівень і конфігурація винесені в окремі модулі.
  \item \textbf{Liskov Substitution}. Реалізації \texttt{PostgresStorage} і \texttt{InMemoryStorage} можуть використовуватися через спільний контракт \texttt{BaseStorage}, а fake-, unavailable- і transformers-сервіси --- через контракт \texttt{BaseLLMService}.
  \item \textbf{Interface Segregation}. Зовнішні рівні системи працюють з обмеженими контрактами, потрібними саме для їхньої задачі: storage-рівень відповідає за дані, а LLM-рівень --- за генерацію.
\end{itemize}

У програмі застосовано кілька практичних архітектурних патернів. Патерн \textbf{Repository/Storage Adapter} використано у рівні збереження даних: доменна логіка працює через \texttt{BaseStorage}, а конкретна реалізація \texttt{PostgresStorage} приховує SQL-запити та структуру таблиць. Патерн \textbf{Service Layer} реалізовано через \texttt{ChatManager}, який координує роботу між інтерфейсом, сховищем і LLM-сервісом. Патерн \textbf{Adapter} також використовується на LLM-рівні: transformers-сервіс приводить зовнішній API Hugging Face до внутрішнього інтерфейсу \texttt{BaseLLMService}.

Принцип \textbf{Inversion of Control} реалізовано через те, що ключові залежності створюються на рівні композиції застосунку, а не всередині доменної логіки. Web-застосунок і CLI-режим отримують готовий \texttt{ChatManager}, а сам \texttt{ChatManager} отримує storage- та LLM-сервіси через конструктор. Завдяки цьому у production-режимі можна використати PostgreSQL- і Transformers-реалізації, а у тестах --- in-memory сховище та fake LLM-сервіс.

Принцип \textbf{Composition over Inheritance} проявляється у тому, що поведінка системи формується через поєднання незалежних компонентів, а не через глибоку ієрархію наслідування. Наприклад, \texttt{ChatManager} не наслідує storage або LLM-класи, а містить посилання на відповідні сервіси. \texttt{LLMRuntime} також не є спеціалізованою моделлю, а керує поточним LLM-сервісом як окремою залежністю.

Принципи \textbf{Loose Coupling} і \textbf{High Cohesion} забезпечуються тим, що кожен модуль має вузьку зону відповідальності та взаємодіє з іншими модулями через стабільні інтерфейси. Це дозволяє окремо змінювати web-інтерфейс, storage-рівень, runtime-керування моделлю або конфігурацію моделей без переписування всієї системи.

Окремо використано принцип конфігураційного розширення: локальні моделі та LoRA-адаптери не жорстко прописані в коді, а описуються через \texttt{model\_config.json} і змінні середовища. Це спрощує зміну моделі без зміни програмної логіки та робить систему зручнішою для подальших експериментів.

\subsection{Реалізація web-інтерфейсу}

Web-інтерфейс реалізовано за допомогою Flask. Функція \texttt{create\_app} створює застосунок, ініціалізує \texttt{ChatManager}, \texttt{LLMRuntime} та реєструє HTTP routes.

Основні маршрути web-застосунку:
\begin{itemize}
  \item \texttt{GET /} --- головна сторінка зі списком чатів;
  \item \texttt{POST /chats} --- створення нового чату;
  \item \texttt{GET /chat/<chat\_id>} --- відкриття конкретного чату;
  \item \texttt{POST /chat/<chat\_id>/send} --- надсилання повідомлення без streaming;
  \item \texttt{POST /chat/<chat\_id>/stream} --- потокова генерація відповіді;
  \item \texttt{POST /api/generation/stop} --- зупинка активної генерації;
  \item \texttt{POST /chat/<chat\_id>/rename} --- перейменування чату;
  \item \texttt{POST} або \texttt{DELETE /chat/<chat\_id>/delete} --- видалення чату;
  \item \texttt{GET /chat/<chat\_id>/export} --- експорт чату у Markdown;
  \item \texttt{GET /api/model/status} --- отримання стану моделі;
  \item \texttt{POST /api/model/switch} --- перемикання моделі;
  \item \texttt{POST /api/model/unload} --- вивантаження моделі з пам'яті;
  \item \texttt{GET /api/messages/search} --- пошук повідомлень;
  \item \texttt{POST /api/generation-preset} --- зміна preset генерації.
\end{itemize}

Клієнтська частина реалізована у файлах \texttt{script.js} та \texttt{style.css}. JavaScript-код відповідає за надсилання повідомлень без перезавантаження сторінки, streaming-відображення відповіді, керування кнопкою зупинки генерації, пошук повідомлень, перейменування чатів, перемикання моделей та показ toast-повідомлень.

У web-інтерфейсі реалізовано відображення Markdown, блоків коду, кнопки копіювання коду та математичних формул. Для рендерингу формул використовується локально підключений KaTeX, що дозволяє відображати LaTeX-подібні формули без звернення до зовнішніх CDN.

\subsection{Реалізація CLI-режиму}

Консольний режим реалізовано у модулі \texttt{src/cli/cli\_app.py}. Він використовує той самий \texttt{ChatManager}, що й web-інтерфейс, тому логіка створення чатів, вибору наявних чатів, збереження повідомлень і генерації відповідей залишається спільною.

CLI-режим надає користувачу просте текстове меню:
\begin{enumerate}
  \item створити новий чат;
  \item вибрати існуючий чат;
  \item завершити роботу.
\end{enumerate}

Після вибору чату користувач може вводити повідомлення у командному рядку. Команда \texttt{/exit} повертає до головного меню, а \texttt{/quit} завершує роботу програми. Наявність CLI-режиму підтверджує, що бізнес-логіка системи не залежить від web-рівня.

\subsection{Реалізація сховища даних}

Для збереження історії діалогів використовується PostgreSQL. Рівень збереження реалізовано через інтерфейс \texttt{BaseStorage} і конкретну реалізацію \texttt{PostgresStorage}. Такий підхід відповідає патерну Repository/Storage Adapter і приховує SQL-запити від доменної логіки.

У базі даних використовуються таблиці \texttt{chats} і \texttt{messages}. Таблиця \texttt{messages} має поле \texttt{position}, яке визначає порядок повідомлень у межах одного чату. Для підвищення ефективності запитів створено індекси за датою оновлення чатів, позицією повідомлень у чаті та timestamp повідомлень.

Схема бази даних керується через Alembic. Під час старту \texttt{PostgresStorage} виконує застосування актуальних міграцій. Це дозволяє уникнути ручного створення таблиць після первинного створення самої бази даних.

\subsection{Реалізація LLM-рівня}

LLM-рівень побудовано навколо абстрактного класу \texttt{BaseLLMService}. Він визначає методи для звичайної та потокової генерації:
\begin{itemize}
  \item \texttt{generate\_response};
  \item \texttt{generate\_response\_stream};
  \item \texttt{finalize\_streamed\_response}.
\end{itemize}

Основною реалізацією є \texttt{TransformersLLMService}. Вона завантажує tokenizer і модель через Hugging Face Transformers, обирає пристрій виконання, формує prompt, запускає генерацію та очищує отриману відповідь. Для instruction-моделей використовується chat template tokenizer-а, а якщо він недоступний, застосовується fallback-побудова plain prompt.

Для випадків, коли модель не завантажилася, використовується fallback LLM-сервіс. Він зберігає інформацію про помилку і дозволяє web-інтерфейсу показати коректний статус замість аварійного завершення застосунку.

Система підтримує generation presets:
\begin{itemize}
  \item \texttt{precise} --- більш сфокусована генерація;
  \item \texttt{balanced} --- збалансований режим;
  \item \texttt{creative} --- більш вільна генерація з більшою кількістю токенів.
\end{itemize}

Також реалізовано підтримку PEFT LoRA/QLoRA-адаптерів. У цьому випадку базова модель завантажується з каталогу \texttt{models/}, а adapter підключається через API бібліотеки PEFT. Налаштування комбінацій базових моделей і адаптерів виконується через конфігураційний файл моделей.

\subsection{Runtime-керування моделлю та потокова генерація}

Оскільки локальна LLM є ресурсом, який може займати значний обсяг оперативної або відеопам'яті, у системі реалізовано окремий компонент \texttt{LLMRuntime}. Він відповідає за стан моделі, перемикання, вивантаження, блокування конфліктних операцій та зупинку активної генерації.

Основні стани runtime:
\begin{itemize}
  \item \texttt{ready} --- модель готова до генерації;
  \item \texttt{loading} --- виконується завантаження або перемикання моделі;
  \item \texttt{error} --- модель не завантажилась через помилку;
  \item \texttt{not\_loaded} --- модель вивантажена.
\end{itemize}

Перемикання моделей виконується асинхронно, щоб web-застосунок не блокувався під час тривалого завантаження. Під час генерації заборонено перемикати або вивантажувати модель, оскільки це могло б призвести до некоректної роботи inference-процесу.

Для покращення користувацького досвіду реалізовано streaming generation. Відповідь моделі передається у браузер частинами через NDJSON-потік. Користувач бачить відповідь поступово, а не лише після завершення генерації. Крім того, реалізовано кнопку зупинки генерації, яка встановлює спеціальний \texttt{Event}, що перевіряється під час роботи Transformers.

Окремо реалізовано background generation jobs. Генерація виконується у фоні, а HTTP stream лише читає події з черги. Це підвищує надійність: навіть якщо клієнтський stream обірветься, backend може завершити генерацію і зберегти відповідь у чат. На рівні інтерфейсу під час активної генерації заблоковано перехід до інших чатів і сторінок, щоб користувач не втрачав контекст поточної відповіді.

\subsection{Конфігурація та локальні моделі}

Конфігурація системи зосереджена у файлі \texttt{src/config.py}. Параметри можуть задаватися через змінні середовища або файл \texttt{.env}. Основні параметри:
\begin{itemize}
  \item \texttt{DATABASE\_URL} або окремі параметри PostgreSQL;
  \item \texttt{LLM\_BACKEND};
  \item \texttt{MODEL\_NAME};
  \item \texttt{ADAPTER\_PATH};
  \item \texttt{GENERATION\_PRESET};
  \item \texttt{MAX\_NEW\_TOKENS}, \texttt{TEMPERATURE}, \texttt{TOP\_P};
  \item \texttt{DEVICE};
  \item \texttt{SYSTEM\_PROMPT}.
\end{itemize}

Web-інтерфейс автоматично сканує каталог \texttt{models/} і показує локально доступні моделі. Окремо файл \texttt{model\_config.json} дозволяє описати наперед підготовлені конфігурації, зокрема поєднання базової моделі та LoRA-адаптера.

\subsection{Тестування програмної системи}

Для перевірки працездатності системи використано автоматизовані тести на базі \texttt{pytest}. Тести охоплюють як доменну логіку, так і web endpoints, storage-рівень, LLM-runtime, побудову prompt-ів, конфігурацію та допоміжні скрипти.

Станом на момент підготовки звіту повний набір тестів містить 105 успішних перевірок. Основні групи тестів наведено у таблиці~\ref{tab:test-groups}.

\begin{longtable}{|p{5cm}|p{10cm}|}
\caption{Групи автоматизованих тестів}
\label{tab:test-groups}\\
\hline
\textbf{Група тестів} & \textbf{Що перевіряється} \\
\hline
\endfirsthead
\hline
\textbf{Група тестів} & \textbf{Що перевіряється} \\
\hline
\endhead
\texttt{chat\_manager} & Створення, перейменування, видалення чатів, додавання повідомлень, streaming-логіка, пошук. \\
\hline
\texttt{web\_app} & Flask endpoints, web-операції, streaming, stop generation, model switching, export, search. \\
\hline
\texttt{llm\_runtime} & Runtime-стани моделі, async switch, unload, stop event, помилки завантаження. \\
\hline
\texttt{transformers\_prompting} & Chat template, fallback prompt, очищення відповіді, identity guard, валідація adapter path. \\
\hline
\texttt{model\_registry} & Сканування локальних моделей і читання \texttt{model\_config.json}. \\
\hline
\texttt{migrations} та \texttt{storage\_factory} & PostgreSQL URL, Alembic upgrade, вимоги до конфігурації сховища. \\
\hline
\texttt{config\_presets} & Generation presets і parsing конфігураційних значень. \\
\hline
\texttt{download\_adapter\_script} & Допоміжний скрипт завантаження PEFT/LoRA-адаптерів. \\
\hline
\end{longtable}
\normalsize

Для ізоляції тестів використовуються \texttt{InMemoryStorage} та \texttt{FakeLLMService}. Це дозволяє перевіряти бізнес-логіку без реальної PostgreSQL-бази та без завантаження важкої локальної моделі. PostgreSQL integration tests винесено в окрему групу і потребують наявності змінної \texttt{DATABASE\_URL}.

Окремо було сформовано набір тестових сценаріїв, які відображають основні функціональні можливості та критичні точки системи. Приклади тестових сценаріїв наведено у таблиці~\ref{tab:test-cases}.

\small
\begin{longtable}{|p{1.2cm}|p{2.7cm}|p{4.2cm}|p{3.9cm}|p{1.3cm}|}
\caption{Тестові сценарії програмної системи}
\label{tab:test-cases}\\
\hline
\textbf{ID} & \textbf{Об'єкт тестування} & \textbf{Сценарій} & \textbf{Очікуваний результат} & \textbf{Ст.} \\
\hline
\endfirsthead
\hline
\textbf{ID} & \textbf{Об'єкт тестування} & \textbf{Сценарій} & \textbf{Очікуваний результат} & \textbf{Ст.} \\
\hline
\endhead
TC-01 & \texttt{ChatManager} & Створення нового чату з коректною назвою. & Чат створено, має унікальний ідентифікатор і доступний у списку чатів. & OK \\
\hline
TC-02 & \texttt{ChatManager} & Спроба створити чат з порожньою або дубльованою назвою. & Система повертає помилку валідації і не створює некоректний запис. & OK \\
\hline
TC-03 & \texttt{ChatManager} & Перейменування існуючого чату. & Назва оновлюється, дата зміни чату оновлюється, дублікати блокуються. & OK \\
\hline
TC-04 & \texttt{ChatManager} & Видалення чату разом з повідомленнями. & Чат видаляється зі сховища і більше не повертається у списку. & OK \\
\hline
TC-05 & \texttt{ChatManager} & Надсилання повідомлення користувача і отримання відповіді асистента. & Обидва повідомлення збережено у правильному порядку. & OK \\
\hline
TC-06 & Web endpoints & Створення, відкриття, перейменування та видалення чату через HTTP routes. & Web-запити повертають коректні статуси і змінюють стан системи. & OK \\
\hline
TC-07 & Web streaming & Надсилання повідомлення через streaming endpoint. & Відповідь повертається частинами, після завершення фінальний текст збережено в історії. & OK \\
\hline
TC-08 & Stop generation & Натискання кнопки зупинки під час активної генерації. & Runtime отримує stop event, генерація завершується контрольовано. & OK \\
\hline
TC-09 & Search & Пошук повідомлень за текстовим запитом. & Система повертає релевантні повідомлення з інформацією про чат і роль автора. & OK \\
\hline
TC-10 & Export & Експорт історії чату у Markdown. & Користувач отримує Markdown-файл з повідомленнями та часовими мітками. & OK \\
\hline
TC-11 & \texttt{LLMRuntime} & Перемикання моделі під час неактивної генерації. & Runtime переходить у стан завантаження, після успіху показує нову активну модель. & OK \\
\hline
TC-12 & \texttt{LLMRuntime} & Спроба перемкнути або вивантажити модель під час генерації. & Операція блокується, щоб не порушити активний inference-процес. & OK \\
\hline
TC-13 & LLM service & Побудова prompt через chat template або fallback-формат. & Prompt формується з урахуванням системного повідомлення та історії діалогу. & OK \\
\hline
TC-14 & Model registry & Сканування каталогу локальних моделей і читання \texttt{model\_config.json}. & У sidebar доступні локальні моделі та описані конфігурації з LoRA-адаптерами. & OK \\
\hline
TC-15 & Storage layer & Робота \texttt{PostgresStorage} і фабрики сховища. & Сховище коректно створює, читає, оновлює і видаляє чати та повідомлення. & OK \\
\hline
TC-16 & Alembic migrations & Перевірка застосування міграцій бази даних. & Схема PostgreSQL приводиться до актуального стану без ручного створення таблиць. & OK \\
\hline
TC-17 & Generation presets & Зміна preset генерації через API. & Параметри генерації оновлюються і використовуються під час наступних відповідей. & OK \\
\hline
TC-18 & Adapter loading & Перевірка конфігурації LoRA/QLoRA-адаптера. & Система валідує шлях до адаптера і підключає його до відповідної базової моделі. & OK \\
\hline
\end{longtable}
\normalsize

Для оцінювання кількісного покриття було використано інструмент \texttt{coverage.py}. Повний запуск тестів виконувався командою:

\begin{lstlisting}[style=plantumlstyle]
python -m coverage run -m pytest
\end{lstlisting}

Після виконання тестів було сформовано звіт покриття для production-коду системи:

\begin{lstlisting}[style=plantumlstyle]
python -m coverage report --include="src/*,migrations/*,scripts/*"
\end{lstlisting}

Результат запуску: \textbf{105 тестів успішно пройдено}, загальне покриття production-коду становить \textbf{77\%}. Це відповідає рекомендованому рівню покриття 70--80\% для основної бізнес-логіки та ключових інтеграційних точок. Найвище покриття мають модулі доменної логіки, конфігурації, runtime-керування моделлю, registry локальних моделей та web-рівень. Нижче покриття має \texttt{TransformersLLMService}, оскільки частина його коду залежить від фактичного завантаження локальної LLM і GPU/CPU inference-середовища, що у unit-тестах ізолюється через fake-сервіси.

\subsection{Порядок запуску та експлуатації}

Для запуску системи необхідно встановити Python 3.10 або вище, створити віртуальне середовище, встановити залежності з \texttt{requirements.txt}, створити PostgreSQL-базу даних і налаштувати файл \texttt{.env}.

Приклад запуску web-режиму:
\begin{lstlisting}[style=plantumlstyle]
python app.py
\end{lstlisting}

Приклад запуску CLI-режиму:
\begin{lstlisting}[style=plantumlstyle]
python app.py -t
\end{lstlisting}

Приклад запуску тестів:
\begin{lstlisting}[style=plantumlstyle]
pytest
\end{lstlisting}

Для локальної роботи з instruction-моделлю рекомендовано використовувати модель Qwen2.5-1.5B-Instruct. Вона розміщується у каталозі локальних моделей \texttt{models/}. Моделі та адаптери є локальними артефактами і не додаються до Git.

Таким чином, реалізована система відповідає поставленим вимогам: підтримує web- і CLI-режими, зберігає історію у PostgreSQL, виконує генерацію через локальний LLM-backend, дозволяє перемикати моделі, використовувати presets та LoRA-адаптери, а також має автоматизоване тестове покриття основних компонентів.

\newpage
